\chapter{Culture and Decision-Making}\label{ch:fm:culture-and-decision-making}

In a two-year geopolitical forecasting tournament in which five university-based research groups competed to produce the most accurate
probabilities, one group reported that three interventions made its forecasters better: training them in probability, letting them share
information and argue their reasons in teams, and putting the best of them together. Each improved both how well the forecasters' probabilities
matched outcomes and how sharply they separated events that happened from those that did not; combined with statistical aggregation, the approach
produced the best forecasts two years in a row. Writing a number down, being scored on it and arguing in the open is most of what a betting culture
is, and a trading firm is a place where every decision is a bet.

\section{Betting cultures: probabilities, not opinions}

A trading desk decides under uncertainty all day: whether a signal is real, whether a strategy is broken (chapter 8), whether to hire, to enter a
market (chapter 24), to raise a limit (chapter 12). A culture that states these beliefs as probabilities can score them later; one that states them
as opinions (``I think it'll work'') cannot, and learns only from outcomes, which mix skill and luck (chapter 10).

\begin{definition}[Decision journal]\label{def:fm:culture-and-decision-making:journal}
A \emph{decision journal}\index{decision journal} records each significant decision when it is made: the options considered, the probabilities
and expected values assigned to their outcomes, the reasons, who decided, and the date; the outcome is added later, so that decisions can be
scored on the information available at the time.
\end{definition}

\section{Scoring forecasts and decisions}

\begin{definition}[Brier score]\label{def:fm:culture-and-decision-making:brier}
The \emph{Brier score}\index{Brier score} of probability forecasts $p_i$ of binary outcomes $y_i\in\{0,1\}$ is the mean squared error
$\frac1N\sum_i(p_i-y_i)^2$: zero for perfect forecasts, 0.25 for a constant forecast of one half, lower is better.
\end{definition}

Murphy's decomposition (1973) splits the score into three parts: the uncertainty of the events themselves, $\bar o(1-\bar o)$, which no forecaster
controls; the reliability, the mean squared gap between the forecasts in each probability bin and the frequency with which those events happened,
which calibration training reduces; and the resolution, the spread of the bins' observed frequencies around the overall one, which knowledge
increases. The score is reliability minus resolution plus uncertainty (\cref{lst:fm:culture-and-decision-making:score}); Book 7's calibration
curve draws the reliability term.

\omcode{code/firm/decisionlog/firm_decisionlog.py}{30}{57}{The Brier score, the calibration bins and Murphy's decomposition.}\label{lst:fm:culture-and-decision-making:score}

\begin{definition}[Forecast aggregation]\label{def:fm:culture-and-decision-making:agg}
\emph{Forecast aggregation}\index{forecast aggregation} combines several forecasters' probabilities for the same event into one: by averaging the
probabilities, by averaging their log-odds, or by averaging the log-odds and multiplying them by a factor above one (extremising) to undo the
moderation that averaging produces when the forecasters hold different information.
\end{definition}

\begin{proposition}[Five opinions or one]\label{prop:fm:culture-and-decision-making:rho}
If $K$ forecasters' errors in log-odds each have variance $\sigma^2$ and pairwise correlation $\rho$, the error of their average has variance
\[
\sigma^2\Bigl(\rho+\frac{1-\rho}{K}\Bigr),
\]
which falls from $\sigma^2$ for one forecaster to $\sigma^2/K$ for independent ones, and stays at $\sigma^2$ whatever $K$ when $\rho=1$.
\end{proposition}

\begin{proof}
$\mathrm{Var}\bigl(\frac1K\sum_ke_k\bigr)=\frac1{K^2}\bigl(K\sigma^2+K(K-1)\rho\sigma^2\bigr)=\sigma^2\bigl(\rho+\frac{1-\rho}K\bigr)$.
\end{proof}

A discussion before forecasting correlates the errors: people converge on the loudest argument, the first number spoken, or the senior person's
view. The value of five opinions is in their independence, which is why good forecasting practice collects the numbers before the discussion and
the discussion before a second round.

\section{Tutorial: five opinions or one}\label{tut:fm:culture-and-decision-making:tut}

\begin{tutorial}
\textbf{Goal.} Score a year of a desk's forecasts, decompose the scores, and compare aggregating five independent forecasts with recording one
consensus after a discussion. \textbf{End state:} calibration curves (\cref{fig:fm:culture-and-decision-making:calib}) and the \omterm{def:fm:culture-and-decision-making:brier}{Brier scores} by rule
(\cref{fig:fm:culture-and-decision-making:rho}).
\begin{tutsteps}
\item \textbf{The events.} 400 events a year, with true probabilities drawn from a Beta(2, 2); the truth itself scores 0.214, and a constant
one half 0.25 (\texttt{fm\_culture.year}, synthetic).
\item \textbf{The forecasters.} Five see the truth's log-odds with noise of standard deviation 0.8, and stretch their log-odds by 1.0 to 1.6: all
but the first are overconfident.
\item \textbf{The scores.} \texttt{firm.decisionlog.brier}, \texttt{murphy} and \texttt{calibration}; \texttt{aggregate} three ways.
\item \textbf{The discussion.} \texttt{fm\_culture.five\_or\_one} correlates the forecasters' errors at $\rho$ from 0 to 1 and scores the consensus
over twenty simulated years.
\end{tutsteps}
\end{tutorial}

\begin{omfigure}
\begin{tikzpicture}
\begin{axis}[width=8cm,height=7cm,xlabel={\small forecast probability},ylabel={\small observed frequency},xmin=0,xmax=1,ymin=0,ymax=1,
  tick label style={font=\footnotesize},grid=major,grid style={gray!20},table/col sep=comma,
  legend cell align=left,legend style={font=\scriptsize,at={(1.05,0.5)},anchor=west}]
\addplot[gray,dashed] coordinates {(0,0) (1,1)};
\addplot[omThm,thick,mark=square*,mark size=1.3pt] table[x=forecast,y=observed]{figdata/desk/26-culture-and-decision-making/calib_f1.csv};
\addplot[omProp,thick,mark=triangle*,mark size=1.6pt] table[x=forecast,y=observed]{figdata/desk/26-culture-and-decision-making/calib_f5.csv};
\addplot[omDef,very thick,mark=*,mark size=1.3pt] table[x=forecast,y=observed]{figdata/desk/26-culture-and-decision-making/calib_agg.csv};
\legend{perfect calibration,forecaster 1,forecaster 5 (overconfident),log-odds average}
\end{axis}
\end{tikzpicture}

\omcaption{Calibration curves for a year of 400 forecasts: forecaster 5, who stretches its log-odds by 1.6, gives high probabilities to events that happen
far less often; the average of five forecasters is close to the diagonal. Data: \texttt{firm.decisionlog.calibration}.}\label{fig:fm:culture-and-decision-making:calib}
\end{omfigure}

In the chapter's year the five forecasters score between 0.225 and 0.262, against 0.214 for the truth itself. The most overconfident, forecaster 5,
has a reliability term of 0.040 against 0.017 for forecaster 1, with about the same resolution (0.026 and 0.027): its extra error is all
miscalibration. Averaging the five probabilities scores 0.220 and averaging their log-odds 0.223, both better than every forecaster
(\cref{fig:fm:culture-and-decision-making:brier}); the average's
reliability term falls to 0.010 and its resolution rises to 0.039, since the independent errors partly cancel. Extremising the average makes it worse
(0.237): these forecasters are already overconfident, and extremising only helps forecasters whose average is too timid.

\begin{omfigure}
\begin{tikzpicture}
\begin{axis}[width=10.5cm,height=5.4cm,xbar,bar width=7pt,xmin=0.20,xmax=0.27,xlabel={\small Brier score (lower is better)},ytick={0,1,2,3,4,5,6,7,8},
  yticklabels={forecaster 1,forecaster 2,forecaster 3,forecaster 4,forecaster 5,mean,log-odds mean,extremised,truth},y dir=reverse,ymin=-0.6,ymax=8.6,
  tick label style={font=\footnotesize},x tick label style={/pgf/number format/.cd,fixed,precision=2},grid=major,grid style={gray!20},
  table/col sep=comma,nodes near coords,nodes near coords style={font=\scriptsize,/pgf/number format/.cd,fixed,precision=3},point meta=x]
\addplot[fill=omDef!55,draw=omDef] table[y=pos,x=brier]{figdata/desk/26-culture-and-decision-making/brier.csv};
\end{axis}
\end{tikzpicture}

\omcaption{\omterm{def:fm:culture-and-decision-making:brier}{Brier scores} of a year's 400 forecasts: five forecasters, three ways of aggregating them, and the truth's own probabilities, the floor
no forecaster can beat on average. Data: \texttt{fm\_culture.scores}.}\label{fig:fm:culture-and-decision-making:brier}
\end{omfigure}

\begin{omfigure}
\begin{tikzpicture}
\begin{axis}[width=10.5cm,height=5cm,xlabel={\small correlation of errors after discussion, $\rho$},ylabel={\small Brier score},xmin=0,xmax=1,ymin=0.20,ymax=0.245,
  tick label style={font=\footnotesize,/pgf/number format/.cd,fixed,precision=3},grid=major,grid style={gray!20},table/col sep=comma,
  legend cell align=left,legend style={font=\scriptsize,at={(0.5,-0.3)},anchor=north,/tikz/every even column/.append style={column sep=8pt}},legend columns=2]
\addplot[omDef,very thick,mark=*,mark size=1.3pt] table[x=rho,y=one]{figdata/desk/26-culture-and-decision-making/rho.csv};
\addplot[omProp,thick,dashed] table[x=rho,y=independent]{figdata/desk/26-culture-and-decision-making/rho.csv};
\legend{one consensus after discussion,five independent forecasts}
\end{axis}
\end{tikzpicture}

\omcaption{The \omterm{def:fm:culture-and-decision-making:brier}{Brier score} of the desk's forecast, averaged over twenty simulated years: five independent forecasts aggregated, against one
consensus whose errors a discussion has correlated at $\rho$. Data: \texttt{fm\_culture.five\_or\_one}.}\label{fig:fm:culture-and-decision-making:rho}
\end{omfigure}

Over twenty simulated years five independent forecasts average a \omterm{def:fm:culture-and-decision-making:brier}{Brier score} of 0.207; a consensus after a discussion that correlates the errors at
0.5 scores 0.224, and at 1 (everyone repeats the same view) 0.237 (\cref{prop:fm:culture-and-decision-making:rho}). The difference, 0.017 at
$\rho=0.5$, is more than half the gap between the average forecaster's score in the chapter's year and the truth's: the discussion throws away most of what the independent
forecasts were worth.

\section{Post-mortems and pre-mortems}

\begin{definition}[Outcome bias]\label{def:fm:culture-and-decision-making:outcome}
\emph{Outcome bias}\index{outcome bias} is judging the quality of a decision by its outcome rather than by the information and reasoning available
when it was made, so that good decisions that lost are called mistakes and bad decisions that won are rewarded.
\end{definition}

Baron and Hershey (1988) studied the bias experimentally; the \omterm{def:fm:culture-and-decision-making:journal}{decision journal} is its antidote. In the chapter's year forecaster 2, the best, bets at
even odds on the 362 events it rates above 55\% or below 45\%: every bet has a positive expected value by its own forecast, and 123 of them, 34\%,
lose. A reviewer who judges by outcome calls a third of the desk's good decisions bad. The same journal shows forecaster 2's real flaw: it wins
66\% of its bets while its forecasts average 77.7\%, overconfidence that only a record of stated probabilities can reveal.

\begin{definition}[Pre-mortem, red team]\label{def:fm:culture-and-decision-making:premortem}
A \emph{pre-mortem}\index{pre-mortem} is a review held before a decision is implemented, in which the team imagines that it has failed and writes
down why, to surface risks that optimism hides. A \emph{red team}\index{red team} is a group assigned to argue against a plan or to attack a
system, independently of the people who made it.
\end{definition}

Post-mortems (Book 15's blameless reviews) look back at an outcome; \omterm{def:fm:culture-and-decision-making:premortem}{pre-mortems} look forward at a decision. Both work only if they separate the
quality of the decision from the luck of the outcome, which requires the journal.

\section{How good firms disagree}

A firm disagrees well when it collects views before discussing them, states them as numbers, lets a junior person's forecast count as much as a
senior's, assigns someone to argue the other side, and scores everyone later. Book 7's research review and pre-registration are the same
discipline applied to research; the research log is a \omterm{def:fm:culture-and-decision-making:journal}{decision journal} for strategies.

\begin{method}[Running a betting culture]\label{met:fm:culture-and-decision-making:run}
\begin{enumerate}
\item Keep a \omterm{def:fm:culture-and-decision-making:journal}{decision journal} for every significant decision, with probabilities and expected values written before the outcome.
\item Collect forecasts independently before meetings; discuss; collect again; aggregate the log-odds.
\item Score forecasters quarterly with the \omterm{def:fm:culture-and-decision-making:brier}{Brier score} and its decomposition; train calibration where the reliability term is large.
\item Review decisions by the process and the information available, not the outcome; hold \omterm{def:fm:culture-and-decision-making:premortem}{pre-mortems} for large decisions and give a \omterm{def:fm:culture-and-decision-making:premortem}{red
team} the right to be heard.
\item Reward good calibration and good process, not lucky outcomes (chapter 10).
\end{enumerate}
\end{method}

\section{Build: the decision journal}\label{bld:fm:culture-and-decision-making:decisionlog}

\begin{build}
\bfield{Purpose} A \omterm{def:fm:culture-and-decision-making:journal}{decision journal} and the scoring of its forecasts: \omterm{def:fm:culture-and-decision-making:brier}{Brier scores}, Murphy's decomposition, calibration, aggregation and an
outcome-bias check.
\bfield{Interface} \texttt{firm.decisionlog}: \texttt{Entry}, \texttt{brier}, \texttt{murphy}, \texttt{calibration}, \texttt{aggregate},
\texttt{logit}, \texttt{expit}, \texttt{outcome\_bias}.
\bfield{Rules} Probabilities are recorded before outcomes; the decomposition is exact when forecasts are constant within bins; aggregation is over
the same events.
\bfield{Acceptance tests} \texttt{code/firm/decisionlog/tests/}: the decomposition's identity on binned forecasts; the three aggregations on hand
numbers; the outcome-bias count.
\bfield{Stretch} Scoring continuous forecasts; weighting forecasters by past skill; a journal of limit and hiring decisions linked to chapter 12's
register.
\end{build}

\begin{omsources}
\item B.~A.~Mellers et al., ``Psychological strategies for winning a geopolitical forecasting tournament'', \emph{Psychological Science} 25(5), 2014.
\item G.~W.~Brier, \emph{Monthly Weather Review} 78(1), 1950; A.~H.~Murphy, \emph{Journal of Applied Meteorology} 12, 1973; V.~A.~Satopaa et al.,
\emph{International Journal of Forecasting}, 2014.
\item J.~Baron and J.~C.~Hershey, ``Outcome bias in decision evaluation'', \emph{Journal of Personality and Social Psychology} 54(4), 1988.
\end{omsources}

\section{Exercises}

\begin{exercise}[$\star$]\label{exo:fm:culture-and-decision-making:1}
Compute the \omterm{def:fm:culture-and-decision-making:brier}{Brier score} of the forecast 0.7 for three events of which two happened. What does a constant 0.5 score?
\end{exercise}

\begin{exercise}[$\star$]\label{exo:fm:culture-and-decision-making:2}
Average the probabilities 0.6 and 0.8, and their log-odds. Which is further from one half?
\end{exercise}

\begin{exercise}[$\star$]\label{exo:fm:culture-and-decision-making:3}
With five forecasters whose errors are independent, by what factor does averaging cut the error variance? With a correlation of 0.5?
\end{exercise}

\begin{exercise}[$\star\star$]\label{exo:fm:culture-and-decision-making:4}
Why does forecaster 5 score worse than forecaster 1 although its resolution is about the same?
\end{exercise}

\begin{exercise}[$\star\star$]\label{exo:fm:culture-and-decision-making:5}
Why does extremising hurt in the chapter's year, and when would it help?
\end{exercise}

\begin{exercise}[$\star\star$]\label{exo:fm:culture-and-decision-making:6}
How should a desk run a meeting to decide whether a strategy is broken, given \cref{prop:fm:culture-and-decision-making:rho}?
\end{exercise}

\begin{exercise}[$\star\star\star$]\label{exo:fm:culture-and-decision-making:7}
\emph{Coding.} Score the journal's bets by outcome and by process. What share of positive-expected-value decisions would an outcome-based review call
mistakes?
\end{exercise}

\begin{exercise}[$\star\star\star$]\label{exo:fm:culture-and-decision-making:8}
\emph{Find the flaw.} ``We made money on that trade, so it was a good decision.''
\end{exercise}

\section{Problem: Five Opinions or One}

\begin{problem}[{Weekend problem --- five opinions or one}]\label{pb:fm:culture-and-decision-making:1}
A head of research wants to replace the desk's weekly consensus meeting with a scored forecasting process.

\medskip\noindent\textbf{Part I --- Principles.}
\begin{enumerate}
\item What did the forecasting tournament of the hook find?
\item Define a \omterm{def:fm:culture-and-decision-making:journal}{decision journal}.
\item Define the \omterm{def:fm:culture-and-decision-making:brier}{Brier score} and give its value for a constant one half.
\item State Murphy's decomposition and what each term measures.
\end{enumerate}

\medskip\noindent\textbf{Part II --- The year.}
\begin{enumerate}[resume]
\item Give the five forecasters' scores and the truth's.
\item Decompose the scores of forecasters 1 and 5 and of the average.
\item Compare the three aggregation rules.
\item Read the calibration curves.
\end{enumerate}

\medskip\noindent\textbf{Part III --- The discussion.}
\begin{enumerate}[resume]
\item Define \omterm{def:fm:culture-and-decision-making:agg}{forecast aggregation}.
\item State and prove \cref{prop:fm:culture-and-decision-making:rho}.
\item Give the \omterm{def:fm:culture-and-decision-making:brier}{Brier scores} of independent forecasts and of a consensus at $\rho=0.5$ and $\rho=1$.
\item Why does discussion correlate errors?
\item How would you run the meeting instead?
\end{enumerate}

\medskip\noindent\textbf{Part IV --- Culture.}
\begin{enumerate}[resume]
\item Define \omterm{def:fm:culture-and-decision-making:outcome}{outcome bias} and give the journal's evidence of it.
\item Define a \omterm{def:fm:culture-and-decision-making:premortem}{pre-mortem} and a \omterm{def:fm:culture-and-decision-making:premortem}{red team}.
\item How should pay reflect forecasting skill?
\item What would you record in a journal for a limit increase?
\item What could go wrong with scoring people?
\item State the \emph{named result}: the Brier-score improvement from aggregating five independent forecasts against one discussed forecast whose
errors are correlated at a given level.
\item In two sentences, write the proposal.
\end{enumerate}
\end{problem}

\section{Interview questions}

\begin{interviewq}[$\star$ \iqroles{researcher}]\label{iq:fm:culture-and-decision-making:1}
What is a \omterm{def:fm:culture-and-decision-making:brier}{Brier score}, and what does a score of 0.25 mean?
\end{interviewq}

\begin{interviewq}[$\star$ \iqroles{trader}]\label{iq:fm:culture-and-decision-making:2}
You lost money on a trade you still think was right. How do you tell?
\end{interviewq}

\begin{interviewq}[$\star\star$ \iqroles{researcher}]\label{iq:fm:culture-and-decision-making:3}
Why average log-odds rather than probabilities?
\end{interviewq}

\begin{interviewq}[$\star\star$ \iqroles{researcher, trader}]\label{iq:fm:culture-and-decision-making:4}
Five people's errors are correlated at 0.5. How much does averaging their forecasts reduce the error variance?
\end{interviewq}

\begin{interviewq}[$\star\star$ \iqroles{risk}]\label{iq:fm:culture-and-decision-making:5}
How would you run a \omterm{def:fm:culture-and-decision-making:premortem}{pre-mortem} for a new market entry?
\end{interviewq}

\begin{interviewq}[$\star\star\star$ \iqroles{researcher}]\label{iq:fm:culture-and-decision-making:6}
Derive Murphy's decomposition of the \omterm{def:fm:culture-and-decision-making:brier}{Brier score}.
\end{interviewq}
